% Generated draft: requires mathematical and source review.
\documentclass[11pt]{article}
\usepackage[margin=25mm]{geometry}
\usepackage{fontspec}
\setmainfont{Latin Modern Roman}
\usepackage{amsmath,amssymb,mathtools,unicode-math}
\setmathfont{Latin Modern Math}
\usepackage{enumitem,needspace,xurl,hyperref}
\setlength{\parindent}{0pt}
\setlength{\parskip}{5pt}
\setlength{\emergencystretch}{3em}
\newcommand{\N}{\mathbb{N}}
\newcommand{\Z}{\mathbb{Z}}
\newcommand{\Q}{\mathbb{Q}}
\newcommand{\R}{\mathbb{R}}
\newcommand{\C}{\mathbb{C}}
\newcommand{\F}{\mathbb{F}}
\newcommand{\A}{\mathbb{A}}
\newcommand{\PP}{\mathbb{P}}
\newcommand{\E}{\mathbb{E}}
\newcommand{\eps}{\varepsilon}
\newcommand{\dd}{\,\mathrm{d}}
\begin{document}
\section*{UnsolvedMath ID: 30001008}
\textbf{Source problem code:} catalogue entry
\subsection{Definitions and mathematical statement}
Consider a plane-wave discontinuous Galerkin method for the homogeneous Helmholtz equation, conventionally written
\[
-\Delta u-\kappa^2u=0\quad\text{in }\Omega,
\]
with wavenumber $\kappa>0$. Homogeneity here concerns the interior equation; the boundary data and boundary conditions are not supplied and are not assumed zero.

For a fixed mesh $\mathcal T_h$, a conventional local approximation space on an element $K$ is
\[
V_p(K)=\operatorname{span}\{e^{i\kappa d_{K,j}\cdot x}:1\leq j\leq m_K(p)\},
\qquad |d_{K,j}|=1.
\]
The global space is broken across element interfaces, with coupling through the chosen DG formulation. In the $p$-version, approximation order or the number of directions $m_K(p)$ increases while the mesh is fixed. The record does not specify a relation between its order parameter $p$ and $m_K(p)$, or how directions are distributed.

Develop and analyze this $p$-version: establish convergence and quantitative error estimates as the approximation order increases, and determine the dependence on wavenumber that governs pollution. A possible analysis would specify suitable error and approximation norms, stability constants and best-approximation estimates before deriving a rate. Neither an algebraic nor an exponential rate is imposed by the question. Pollution analysis must distinguish approximation error from deterioration of the numerical method as $\kappa$ grows and state the accompanying resolution regime; a fixed-wavenumber convergence statement alone does not specify that behavior.

The domain and dimension, boundary-value problem and its well-posedness, mesh regularity, numerical fluxes and penalty parameters, direction choices, solution regularity, error norm and joint dependence of $h,p,\kappa$ are all unstated. Thus no particular DG scheme, stability theorem or rate can be uniquely formulated from this record alone.

The supplied assessment checked 2026-08-21 labels the problem partially solved. It describes the historical 2009 question as substantially resolved in important settings by $p$-version analysis and later exponential $hp$ convergence, while saying no single result settles the broad statement as written. That reported $hp$ progress is retained separately from the requested fixed-mesh $p$-version analysis.
\subsection{Short English statement}
Analyze convergence, error rates and pollution when increasing plane-wave approximation order in a fixed-mesh DG Helmholtz method. The source leaves the boundary problem, DG formulation, direction strategy and parameter regime unspecified and reports substantial progress in particular settings.
\subsection{Sources}
\begin{itemize}
\item Local UnsolvedMath registry, permanent ID 30001008, source code catalogue entry; supplied complete record, including its corrected-verified editorial-cleanup note. Literature checked 2026-08-21: the assessment is supplied dated curation, not independently verified here.
\item Nonstandard Finite Element Methods (2009), Oberwolfach Reports, DOI 10.4171/owr/2008/36, pp. 2027--2094: \url{https://doi.org/10.4171/owr/2008/36}. Proposed-year metadata gives 2008.
\item PWDG $p$-version analysis: \url{https://doi.org/10.1137/090761057}.
\item PWDG $hp$ convergence: \url{https://centaur.reading.ac.uk/40203/1/HiptmairMoiolaPerugia_FoCM_accepted.pdf}.
\end{itemize}
\end{document}
